\documentclass[a4paper]{article}

\usepackage{anyfontsize}
\usepackage{amsmath}

\usepackage[no-math]{fontspec}
\setmainfont{XITS}
\usepackage{unicode-math}
\unimathsetup{bold-style=ISO,sans-style=italic}
\setmathfont{XITS Math}
\AtBeginDocument{
  \let\uglymathbf\mathbf
  \renewcommand\mathbf\symbf
  \let\uglymathsf\mathsf
  \renewcommand\mathsf\symsf
}

\AtBeginDocument{
  \let\uglyepsilon\epsilon
  \renewcommand\epsilon\varepsilon
}

\usepackage{bm}
\renewcommand\bm\symbf

\usepackage{indentfirst}
\setlength{\parindent}{0em}
\usepackage{autobreak}
\allowdisplaybreaks
\usepackage{parskip}
\usepackage{geometry}
\geometry{top=3cm,bottom=3cm,left=2cm,right=2cm}

\usepackage[fontset=none]{ctex}
\setCJKmainfont{Adobe Song Std}[BoldFont=Noto Sans CJK SC Medium]

\begin{document}

\title{\textbf{实测天体物理作业}}
\author{1900011604\ \ 张植竣\ \ 北京大学}
\date{2022年3月21日}
\maketitle

\begin{enumerate}

    \item[1.]
    对于圆形孔径：
    \[
        \mathrm{FWHP} = 58.4^\circ \times \frac{\lambda}{D} = 58.4^\circ \times \frac{c}{fD} = 1.75^\circ
    \]

    \item[2.]
    由源张角$\theta_S = 30'$，主瓣张角$\theta_B = 1.75^\circ = 105'$得主瓣温度：
    \[
        T_B = T\left(\frac{\theta_S^2}{\theta_S^2 + \theta_B^2}\right) = 5800\,\mathrm{K} \times \left(\frac{30^2}{30^2 + 105^2}\right) = 437.7\,\mathrm{K}
    \]
    天线温度：
    \[
        T_A = \eta_B T_B = 0.7 \times 437.7\,\mathrm{K} = 306.4\,\mathrm{K}
    \]

    \item[3.]
    条纹间距
    \[
        \delta = \frac{\lambda}{b} = \frac{c}{fb} = 61.9''
    \]
    而复可见度函数：
    \[
        V(u) = \int_{-\tfrac{\theta_S}{2}}^{\tfrac{\theta_S}{2}} A I_0 \mathrm{e}^{2\pi iu\theta}\,\mathrm{d}\theta,\quad u = \frac{b}{\lambda}
    \]
    由于$A$、$I_0$均为常量，代入积分得：
    \[
        V(u) = AI_0\frac{\sin(\pi u\theta_S)}{\pi u} = -3.2\times10^{5}AI_0
    \]
\end{enumerate}

\end{document}